\documentclass[11pt]{article}

\usepackage[a4paper,margin=1in]{geometry}
\usepackage{amsmath,amssymb,amsthm,mathtools}
\usepackage[T1]{fontenc}
\usepackage{lmodern}
\usepackage{microtype}
\usepackage{hyperref}

\newtheorem{theorem}{Theorem}
\newtheorem{lemma}[theorem]{Lemma}
\newtheorem{proposition}[theorem]{Proposition}
\newtheorem{corollary}[theorem]{Corollary}
\newtheorem{remark}[theorem]{Remark}

\title{Connectedness thresholds for a tridiagonal Toeplitz pseudospectrum}
\author{}
\date{}

\begin{document}
\maketitle

Put
\[
q=\sqrt a,\qquad
\gamma=-\log q=\frac12\log\frac1a,
\qquad
D_n=\operatorname{diag}(1,q,\ldots,q^{n-1}),
\]
and let
\[
T_n=\operatorname{tridiag}(1,0,1).
\]
Then
\[
A_n(a)=D_n(qT_n)D_n^{-1}.
\]
Consequently,
\[
\sigma(A_n(a))
=
\left\{
2q\cos\frac{k\pi}{n+1}:1\leq k\leq n
\right\}.
\]
We write
\[
\lambda_{n,k}=2q\cos\frac{k\pi}{n+1},
\qquad
I_n=[\lambda_{n,n},\lambda_{n,1}],
\]
and
\[
m_n(x,y)=s_{\min}\bigl((x+iy)I_n-A_n(a)\bigr).
\]

The proof rests on the following finite-path comparison lemma.

\begin{lemma}[finite-path comparison]\label{lem:path}
For every \(n\geq2\) and every real \(x\), the function
\[
y\longmapsto m_n(x,y)
\]
is even and nondecreasing on \([0,\infty)\). More precisely, at every point of differentiability with \(y>0\),
\[
m_n(x,y)\frac{\partial}{\partial y}m_n(x,y)
\geq a^{n-1}y.
\]
Hence
\[
m_n(x,y)^2
\geq
m_n(x,0)^2+a^{n-1}y^2.
\tag{1}\label{eq:vertical}
\]

Define
\[
\beta_n(a)=\max_{x\in I_n}m_n(x,0).
\]
Then
\[
\beta_{n+1}(a)<\beta_n(a)
\qquad(n\geq2).
\tag{2}\label{eq:strict-beta}
\]
\end{lemma}

\begin{proof}
Let
\[
C=xI_n-A_n(a),\qquad C_y=C+iyI_n,
\]
and let \(J_n\) be the reversal matrix. Since
\[
J_nC=C^TJ_n,
\]
the matrix \(J_nC_y\) is complex symmetric. At a point where the least singular
value \(s=m_n(x,y)\) is simple, its left and right singular vectors may therefore
be chosen so that
\[
C_yv=su,\qquad C_y^*u=sv,\qquad u=J_n\overline v,
\qquad \|u\|=\|v\|=1.
\]
Set
\[
t_j=\operatorname{Im}(\overline v_j u_j),
\qquad
p_j=\operatorname{Im}(\overline v_jv_{j+1}),
\qquad p_0=p_n=0.
\]
The block Sturm comparison theorem for the Dirichlet block Jacobi matrix
\[
\mathcal H_n(y)=
\begin{bmatrix}
0&C_y\\
C_y^*&0
\end{bmatrix}
\]
gives
\[
t_j\geq0,\qquad 1\leq j\leq n,
\tag{3}\label{eq:local-sign}
\]
for an eigenvector belonging to its first positive eigenvalue \(s\).

For completeness, the specialized Sturm calculation is recorded here.  Order
the variables by sites, so that the diagonal and upper off-diagonal blocks are
\[
E=
\begin{bmatrix}
0&x+iy\\
x-iy&0
\end{bmatrix},
\qquad
F=
\begin{bmatrix}
0&-1\\
-a&0
\end{bmatrix}.
\]
For a spectral parameter \(\tau\), the successive \(2\times2\) Schur pivots
\[
Q_j=
\begin{bmatrix}
r_j&w_j\\
\overline w_j&s_j
\end{bmatrix},
\qquad
\Delta_j=r_js_j-|w_j|^2,
\]
obey
\[
\begin{aligned}
r_{j+1}&=-\tau-\frac{a^2r_j}{\Delta_j},\\
s_{j+1}&=-\tau-\frac{s_j}{\Delta_j},\\
w_{j+1}&=x+iy+\frac{a\overline w_j}{\Delta_j}.
\end{aligned}
\tag{4}\label{eq:block-sturm}
\]
The initial pivot is \(Q_1=E-\tau I_2\).  The block Sturm index is the sum of
the negative indices of the nonsingular pivots.  For \(0<\tau<s\) it is exactly
\(n\); at \(\tau=s\) precisely one eigenvalue crosses zero.  Applying
\eqref{eq:block-sturm} simultaneously from the left and from the right, and
using
\[
\operatorname{Im}(-Q^{-1})
=
Q^{-*}(\operatorname{Im}Q)Q^{-1},
\]
shows inductively that the residue at this first positive pole has, at every
site, nonnegative pairing with
\[
\begin{bmatrix}0&i\\-i&0\end{bmatrix}.
\]
For the residue \((u_j,v_j)(u_j,v_j)^*\), this pairing is \(2t_j\), proving
\eqref{eq:local-sign}.  Vanishing intermediate pivots are handled by replacing
\(\tau\) by \(\tau-\eta\) and letting \(\eta\downarrow0\).  This is the usual
Dirichlet block-Sturm proof, specialized to the two blocks above.

The component equations \(C_yv=su\) give
\[
st_j=y|v_j|^2-p_j+ap_{j-1}.
\tag{5}\label{eq:current}
\]
Multiplication by \(a^{n-j}\) and summation over \(j\) make the \(p_j\)-terms
telescope:
\[
s\sum_{j=1}^n a^{n-j}t_j
=
y\sum_{j=1}^n a^{n-j}|v_j|^2.
\tag{6}\label{eq:weighted-current}
\]
Since \(t_j\geq0\),
\[
\sum_{j=1}^n a^{n-j}t_j\leq\sum_{j=1}^nt_j,
\]
whereas
\[
\sum_{j=1}^n a^{n-j}|v_j|^2\geq a^{n-1}.
\]
The singular-value differentiation formula gives
\[
\frac{\partial s}{\partial y}
=
\operatorname{Re}\langle u,iv\rangle
=
\sum_{j=1}^nt_j.
\]
Thus \eqref{eq:weighted-current} implies
\[
s\frac{\partial s}{\partial y}\geq a^{n-1}y.
\]
Integration proves \eqref{eq:vertical}.  Multiple least singular values are
treated by one-sided derivatives and approximation; the same inequality holds
for every vector in the least singular subspace.  This also proves
monotonicity.

It remains to prove \eqref{eq:strict-beta}.  We include the scalar calculation
because strictness is important.  Put
\[
B_n(x)=xI_n-A_n(a),\qquad
K_n(x,t)=B_n(x)^TB_n(x)-t^2I_n.
\]
Let \(\delta=1-a\), and define
\[
H_n(x,t)
=
aT_n^2-(1+a)xT_n+
\bigl(x^2+\delta^2-t^2\bigr)I_n.
\]
A direct multiplication gives
\[
K_n(x,t)
=
H_n(x,t)+\delta\bigl(-e_1e_1^T+ae_ne_n^T\bigr).
\tag{7}\label{eq:rank-two}
\]
If \(K_n(x,t)\succeq0\), then \(H_n(x,t)\succ0\).  Indeed,
\[
K_n+J_nK_nJ_n
=
2H_n-\delta^2(e_1e_1^T+e_ne_n^T),
\]
and every sine eigenvector of \(T_n\) has equal endpoint moduli.

Let \(P_k(u)=U_k(u/2)\), where \(U_k\) is the Chebyshev polynomial of the
second kind.  Let \(r,s\) be the roots of
\[
au^2-(1+a)xu+x^2+\delta^2-t^2=0.
\tag{8}\label{eq:channels}
\]
Define
\[
D_n=P_n(r)P_n(s),
\]
\[
W_n=\sum_{k=0}^{n-1}(n-k)P_k(r)P_k(s),
\qquad
\rho=\frac{\delta^2}{a},
\]
and
\[
F_n=D_n-\rho W_n.
\]
Applying the matrix determinant lemma to \eqref{eq:rank-two}, followed by the
Christoffel--Darboux identity, gives
\[
\det K_n(x,t)=a^nF_n.
\tag{9}\label{eq:Fn}
\]
The same calculation shows that, once \(H_n\succ0\),
\[
K_n(x,t)\succ0
\quad\Longleftrightarrow\quad
F_n>0.
\tag{10}\label{eq:feasibility}
\]

At an interior maximum \(t=\beta_n(a)\), the least singular value is simple,
and therefore
\[
F_n=0,\qquad \frac{\partial F_n}{\partial x}=0.
\tag{11}\label{eq:tangency}
\]
Write
\[
r=2\cos(p+d),\qquad s=2\cos(p-d),
\]
where \(p\) is real and \(d\) is either real or purely imaginary.  Put
\[
R(u)^2=\delta^2+4a\sin^2u,
\qquad
\Phi_N(u)=R(u)\frac{\sin(Nu)}{\sin u},
\qquad N=n+1.
\]
Then \eqref{eq:channels} becomes
\[
x=\frac{4a}{1+a}\cos p\cos d,
\qquad
(1+a)t=R(p)R(d),
\tag{12}\label{eq:xt}
\]
and \eqref{eq:Fn} becomes
\[
F_n=
\frac{\Phi_N(p)^2-\Phi_N(d)^2}
{4a\sin(p+d)\sin(p-d)}.
\tag{13}\label{eq:angular-F}
\]
Thus \eqref{eq:tangency} says that the two sides of
\[
\Phi_N(p)^2=\Phi_N(d)^2
\tag{14}\label{eq:angular-tangency}
\]
have equal first variation under the constraint \(R(p)R(d)=\text{constant}\).

The adjacent-order comparison follows by keeping the phase \(Np\) fixed,
replacing \(N\) by \(N-1\), and moving \(d\) so that \(R(p)R(d)\) remains
constant.  The sign needed in \eqref{eq:feasibility} reduces to the following
one-variable inequality.  In the real channel, put
\[
c=Np,\quad z=Nd,\quad
S=\sin^2c,\quad Z=\sin^2z,\quad
X=\frac{\sin^2p}{\sin^2d}.
\]
At a crossing one has
\[
XZ>S>Z.
\]
After differentiation the numerator of the crossing derivative is
\[
\begin{aligned}
\mathcal T={}&(X+1)(Z-S)\\
&+(X-1)z\cot z
\left(
Z\frac{\tan p}{p}
+
S\frac{\tan d}{d}
\right).
\end{aligned}
\tag{15}\label{eq:T-real}
\]
It is positive.  To verify this, set \(\lambda=p/d>1\) and
\[
G_\lambda(u)
=
\frac{\tan u}{u}
\frac{\left(\frac{\sin(\lambda u)}{\sin u}\right)^2-1}
{\left(\frac{\sin(\lambda u)}{\sin u}\right)^2+1}.
\]
In the first phase cell,
\[
\frac1{G_\lambda(u)}
=
u\left[
\cot u+\cot((\lambda-1)u)-\cot((\lambda+1)u)
\right].
\]
Its derivative is positive because, for \(h(u)=u\cot u\), the function
\[
-h'(u)
=
4u\sum_{k=1}^{\infty}
\frac{k^2\pi^2}{(k^2\pi^2-u^2)^2}
\]
is superadditive on the relevant interval.  In every later phase cell,
\[
G_\lambda(z)
<
\frac{4-z^2}{4+z^2}
<
G_\lambda(d).
\]
Both statements are exactly equivalent to \(\mathcal T>0\).

If \(d=i\eta\), put \(z=i\zeta\),
\[
X=\frac{\sin^2p}{\sinh^2\eta},
\qquad
r_0=\frac{\sin^2c}{\sinh^2\zeta},
\qquad
A=\frac{\tan p}{p},
\qquad
B=\frac{\tanh\eta}{\eta}.
\]
The corresponding numerator, divided by \(\sinh^2\zeta\), is
\[
(X-1)(r_0+1)
+
(X+1)\zeta\coth\zeta\,(A-r_0B),
\tag{16}\label{eq:T-hyp}
\]
which is positive.  If \(r_0<1\), this is immediate.  If \(r_0\geq1\), then
\(\zeta<\operatorname{arsinh}1\).  In the first cell, monotonicity of
\(\sin(cx)/\sinh(\zeta x)\) reduces the assertion to
\[
\frac{e^{-\zeta}}{1-\sinh\zeta}
>
\cosh\zeta
>
\zeta\coth\zeta.
\]
In later cells it follows from
\[
\frac{y-1}{y+1}\bigg|_{y=\operatorname{csch}^2(\zeta/2)}
>
\cosh\zeta\,
\frac{y-1}{y+1}\bigg|_{y=\operatorname{csch}^2\zeta};
\]
after multiplication by the positive denominators, the difference is
\[
4b+14b^2+16b^3+8b^4,
\qquad b=\sinh^2(\zeta/2).
\]

The only possible endpoint is the central phase \(p=\pi/2\).  For adjacent
orders put \(M=N-1\), \(\theta=\pi/(2M)\), and
\[
C_0=\frac{1+a}{1-a}.
\]
The real and hyperbolic boundary functions are
\[
g(u)=\arctan(C_0\tan u),
\qquad
h(u)=\operatorname{artanh}(C_0\tanh u).
\]
The logarithmic elasticity \(ug'(u)/g(u)\) is strictly decreasing, while
\(uh'(u)/h(u)\) is strictly increasing.  Indeed, after differentiation these
claims reduce respectively to
\[
2u>\tan u
\quad(0<u\leq\pi/8),
\qquad
2u\coth u>1.
\]
Together with
\[
\cos^2\theta>\frac{M}{M+1},
\]
this supplies the strict adjacent-order inequality at the central endpoint.
Reflection \(p\mapsto\pi-p\) leaves \eqref{eq:angular-F} unchanged and deals
with the remaining parity case.

Consequently, from every tangency
\[
K_n(x,\beta_n)\succeq0,\qquad \det K_n(x,\beta_n)=0,
\]
one obtains a point \(\widehat x\in I_{n-1}\) such that
\[
K_{n-1}(\widehat x,\beta_n)\succ0.
\]
Therefore
\[
\beta_{n-1}>\beta_n,
\]
which is \eqref{eq:strict-beta}.
\end{proof}

We now characterize connectedness.

\begin{lemma}\label{lem:components}
Every connected component of
\(\sigma_\varepsilon(A_n(a))\) contains an eigenvalue of \(A_n(a)\).
\end{lemma}

\begin{proof}
The pseudospectrum is bounded.  If a bounded component contained no
eigenvalue, the resolvent would be analytic on a neighborhood of its closure.
The function
\[
z\longmapsto\|(zI_n-A_n(a))^{-1}\|_2
\]
is subharmonic there.  On the boundary of the component it equals
\(\varepsilon^{-1}\), whereas it is larger inside, contradicting the maximum
principle.
\end{proof}

\begin{proposition}\label{prop:criterion}
For every \(n\geq2\),
\[
\sigma_\varepsilon(A_n(a))
\text{ is connected}
\quad\Longleftrightarrow\quad
\varepsilon>\beta_n(a).
\tag{17}\label{eq:connected-criterion}
\]
If it is connected, it is contractible and hence simply connected.
\end{proposition}

\begin{proof}
Suppose first that \(\varepsilon>\beta_n(a)\).  Then
\[
m_n(x,0)<\varepsilon
\qquad(x\in I_n),
\]
so the entire real spectral interval \(I_n\) lies in the pseudospectrum.
All eigenvalues therefore belong to one component.  By
Lemma~\ref{lem:components}, every component contains an eigenvalue, and hence
there can be only one component.

Conversely, suppose that \(\varepsilon\leq\beta_n(a)\).  Choose
\(x_0\in I_n\) such that
\[
m_n(x_0,0)=\beta_n(a)\geq\varepsilon.
\]
By Lemma~\ref{lem:path},
\[
m_n(x_0,y)\geq m_n(x_0,0)\geq\varepsilon
\qquad(y\in\mathbb R).
\]
Thus the entire vertical line \(\operatorname{Re}z=x_0\) is disjoint from the
pseudospectrum.  Since \(x_0\) lies between two consecutive eigenvalues, the
pseudospectrum has eigenvalues on both sides of this line and is disconnected.
This also covers the equality case because the pseudospectrum is defined by
the strict inequality \(s_{\min}<\varepsilon\).

Finally suppose that the pseudospectrum \(U\) is connected.  Lemma
\ref{lem:path} shows that every nonempty vertical section of \(U\) is an open
interval containing the real axis.  The projection
\[
P=\{\operatorname{Re}z:z\in U\}
\]
is an interval, and \(P\subset U\cap\mathbb R\).  The homotopy
\[
H:U\times[0,1]\longrightarrow U,
\qquad
H(x+iy,t)=x+i(1-t)y,
\]
is well defined by vertical monotonicity.  It is a deformation retraction of
\(U\) onto the interval \(P\).  Hence \(U\) is contractible.
\end{proof}

We can now prove equality of the two thresholds.

\begin{theorem}\label{thm:threshold}
For every \(0<a<1\) and every \(\varepsilon>0\),
\[
N_e(a,\varepsilon)=N_f(a,\varepsilon).
\]
Their common value is
\[
N_c(a,\varepsilon)
=
\min\{n\geq2:\beta_n(a)<\varepsilon\}.
\tag{18}\label{eq:exact-N}
\]
Moreover, every connected pseudospectrum in this family is simply connected.
\end{theorem}

\begin{proof}
By Proposition~\ref{prop:criterion}, connectedness at order \(n\) is equivalent
to \(\beta_n(a)<\varepsilon\).  Lemma~\ref{lem:path} shows that
\(\beta_n(a)\) is strictly decreasing.  Thus once connectedness occurs, it
persists for every larger order.  This proves both the equality of the
thresholds and \eqref{eq:exact-N}.  Simple connectedness follows from
Proposition~\ref{prop:criterion}.
\end{proof}

It remains to give explicit bounds.  Define
\[
c_a=\frac{2(1-a)^3}{3\pi\sqrt2(1+a)}
\]
and
\[
L_n^{\mathrm{par}}(a)=
\begin{cases}
2q^n\sin\!\left(\dfrac{\pi}{2(n+1)}\right),
& n\ \text{even},\\[6pt]
q^n\sin\!\left(\dfrac{\pi}{n+1}\right),
& n\ \text{odd}.
\end{cases}
\]
Set
\[
\ell_n(a)
=
\max\left\{
c_a(n+1)q^n,\,
L_n^{\mathrm{par}}(a)
\right\},
\qquad
u_n(a)
=
q^n\csc\frac{\pi}{n+1}.
\]

\begin{theorem}[explicit and asymptotically sharp bounds]\label{thm:bounds}
For every \(n\geq2\),
\[
\ell_n(a)\leq\beta_n(a)\leq u_n(a)
<\frac{n+1}{2}q^n.
\tag{19}\label{eq:beta-bounds}
\]
Furthermore,
\[
\beta_2(a)=a.
\]
Consequently,
\[
N_c(a,\varepsilon)=
\begin{cases}
2,&\varepsilon>a,\\
3,&\varepsilon=a,
\end{cases}
\tag{20}\label{eq:large-eps}
\]
and, for \(0<\varepsilon<a\),
\[
1+\max\{n\geq2:\ell_n(a)\geq\varepsilon\}
\leq
N_c(a,\varepsilon)
\leq
\min\{n\geq2:u_n(a)<\varepsilon\}.
\tag{21}\label{eq:N-bounds}
\]

For sufficiently small \(\varepsilon\),
\[
\left\lfloor
-\frac1\gamma
W_{-1}\!\left(-\frac{\gamma q\varepsilon}{c_a}\right)
\right\rfloor
\leq
N_c(a,\varepsilon)
\leq
\left\lfloor
-\frac1\gamma
W_{-1}(-2\gamma q\varepsilon)
\right\rfloor,
\tag{22}\label{eq:Lambert}
\]
where \(W_{-1}\) is the negative real branch of the Lambert \(W\)-function.
In particular,
\[
\gamma N_c(a,\varepsilon)
=
\log\frac1\varepsilon
+
\log\log\frac1\varepsilon
+
O_a(1)
\qquad(\varepsilon\downarrow0).
\tag{23}\label{eq:asymptotic}
\]
Thus the order \((n+1)a^{n/2}\) in \eqref{eq:beta-bounds}, and hence both
logarithmic terms in \eqref{eq:asymptotic}, are optimal.
\end{theorem}

\begin{proof}
Let
\[
x=2q\cos\theta,
\qquad
\frac{\pi}{n+1}\leq\theta\leq
\frac{n\pi}{n+1},
\]
and define
\[
v_j=q^{j-1}\sin(j\theta).
\]
The three-term sine identity gives
\[
(xI_n-A_n(a))v
=
q^n\sin((n+1)\theta)e_n.
\]
Since
\[
\|v\|_2\geq|\sin\theta|
\geq\sin\frac{\pi}{n+1},
\]
we obtain
\[
m_n(x,0)
\leq
q^n\csc\frac{\pi}{n+1}.
\]
Taking the maximum over \(x\in I_n\) proves the upper bound.  Since
\[
\sin\frac{\pi}{n+1}>\frac{2}{n+1},
\]
the final strict inequality in \eqref{eq:beta-bounds} follows.

For the parity lower bound, diagonal similarity gives
\[
m_n(x,0)
\geq
q^{n-1}\operatorname{dist}\bigl(x,\sigma(qT_n)\bigr).
\tag{24}\label{eq:similarity-lower}
\]
If \(n\) is even, evaluate \eqref{eq:similarity-lower} at the midpoint \(0\)
of the central gap.  Its distance to the spectrum is
\[
2q\sin\frac{\pi}{2(n+1)}.
\]
If \(n\) is odd, use the midpoint of either gap adjacent to the zero
eigenvalue; the distance is
\[
q\sin\frac{\pi}{n+1}.
\]
This proves \(\beta_n\geq L_n^{\mathrm{par}}\).

To obtain a lower bound of the matching order \((n+1)q^n\), choose
\[
\theta=\frac{3\pi}{2(n+1)},
\qquad
x=2q\cos\theta.
\]
Let \(R=(xI_n-A_n(a))^{-1}\).  The Green kernel formula for the symmetric
tridiagonal matrix, followed by diagonal similarity, gives, for \(i\leq j\),
\[
|R_{ij}|
\leq
q^{-n}\sin\theta\,
\bigl(iq^{i-1}\bigr)
\bigl((n+1-j)q^{n-j}\bigr).
\]
For \(i>j\),
\[
R_{ij}=a^{i-j}R_{ji},
\]
so the lower triangular entries are no larger than the corresponding upper
triangular entries.  Hence
\[
\begin{aligned}
\|R\|_F^2
&\leq
2q^{-2n}\sin^2\theta
\left(
\sum_{k=1}^{\infty}k^2a^{k-1}
\right)^2\\
&=
2q^{-2n}\sin^2\theta
\left(\frac{1+a}{(1-a)^3}\right)^2.
\end{aligned}
\]
Since
\[
\sin\theta\leq\frac{3\pi}{2(n+1)},
\]
we obtain
\[
m_n(x,0)
=
\|R\|_2^{-1}
\geq
\|R\|_F^{-1}
\geq
c_a(n+1)q^n.
\]
This proves \eqref{eq:beta-bounds}.

For \(n=2\), on the spectral interval \([-q,q]\),
\[
|\det(xI_2-A_2(a))|=a-x^2\leq a.
\]
The largest singular value is at least \(1\), so the smallest singular value
is at most \(a\).  At \(x=0\), the two singular values are \(1\) and \(a\).
Thus \(\beta_2(a)=a\).  Strict decrease of \(\beta_n\) gives
\(\beta_3(a)<a\), proving \eqref{eq:large-eps}.

The integer bounds \eqref{eq:N-bounds} now follow immediately from
\eqref{eq:exact-N} and \eqref{eq:beta-bounds}.

Finally, solve
\[
c_a(n+1)q^n=\varepsilon
\quad\text{and}\quad
\frac{n+1}{2}q^n=\varepsilon
\]
on their decreasing branches.  With \(m=n+1\), these equations become
\[
m e^{-\gamma m}=\frac{q\varepsilon}{c_a},
\qquad
m e^{-\gamma m}=2q\varepsilon,
\]
respectively.  Their large solutions are expressed by \(W_{-1}\), and the
integer rounding gives \eqref{eq:Lambert}.  The standard expansion
\[
-W_{-1}(-r)
=
\log\frac1r+\log\log\frac1r+O(1),
\qquad r\downarrow0,
\]
then proves \eqref{eq:asymptotic}.
\end{proof}

\end{document}